% GENERATED FILE. Edit canonical lessons, then run npm run generate:books.
% canonical-source-hash: fnv1a64:d102dd66e2dddb80
% canonical-lessons: RU-C162-bluzka, RU-C162-futbolka, RU-C162-kolgotki, RU-C162-krossovki, RU-C162-tapochki

\chapter{Blouse, T-shirt, Tights, Trainers, Slippers}
\label{ch:ru-blouse-t-shirt-tights-trainers-slippers}
\hlchaptermodality{162}

\begin{chapteropening}
\textbf{By the end of this chapter:} \emph{I can say a blouse, a T-shirt, tights, trainers and slippers.}

Complete the last lesson of chapter 162: I can say a blouse, a T-shirt, tights, trainers and slippers.
\end{chapteropening}

\section[\texorpdfstring{\ru{блузка} (blúzka)}{блузка (blúzka)}]{\ru{блузка} (blúzka) — a blouse}

\label{lesson:RU-C162-bluzka}



\begin{quote}
Before the new one: say the Russian for a postcard, then the Russian for a parcel.
\end{quote}

\subsection*{You'll want to know: \ru{блузка}}
\textbf{\ru{блузка}} (\emph{blúzka}) — \textquotedblleft{}a blouse\textquotedblright{}.

\begin{cousinweb}
\textbf{\ru{блузка}} is \emph{\ru{блуза}}, \textquotedblleft{}a blouse\textquotedblright{}, with \emph{-\ru{ка}}. \emph{\ru{блуза}} came from German \emph{Bluse} or French \emph{blouse}, which Vasmer, with Kluge, traces, probably, to Medieval Latin \emph{pelusia}, cloth from the Egyptian town of Pelusium.
\end{cousinweb}

\begin{grammarlens}[title={What you've built}]
One more word for this chapter.
\end{grammarlens}

\subsection*{Your turn}
\begin{itemize}
\raggedright
  \item \emph{Say it:} \emph{blúzka}
  \item \emph{Say it:} \emph{blúzka}, once more
  \item \emph{Say it:} say \emph{blúzka}
  \item \emph{Recall:} say the Russian for a postcard, then the Russian for a parcel, then say \emph{blúzka} again
\end{itemize}

\begin{tcolorbox}[breakable,colback=teal!4,colframe=teal!35!black,title={Before you move on}]
What does \ru{блузка} mean? (\textquotedblleft{}A blouse\textquotedblright{}.) Say it once more.
\end{tcolorbox}
\section[\texorpdfstring{\ru{футболка} (futbólka)}{футболка (futbólka)}]{\ru{футболка} (futbólka) — a T-shirt}

\label{lesson:RU-C162-futbolka}



\begin{quote}
Before the new one: say the Russian for a parcel, then the Russian for a blouse.
\end{quote}

\subsection*{You'll want to know: \ru{футболка}}
\textbf{\ru{футболка}} (\emph{futbólka}) — \textquotedblleft{}a T-shirt\textquotedblright{}.

\begin{cousinweb}
\textbf{\ru{футболка}} is \emph{\ru{футбол}}, \textquotedblleft{}football\textquotedblright{}, from English, with \emph{-\ru{ка}}: first the shirt footballers wore.
\end{cousinweb}

\begin{grammarlens}[title={What you've built}]
One more word for this chapter.
\end{grammarlens}

\subsection*{Your turn}
\begin{itemize}
\raggedright
  \item \emph{Say it:} \emph{futbólka}
  \item \emph{Say it:} \emph{futbólka}, once more
  \item \emph{Say it:} say \emph{futbólka}
  \item \emph{Recall:} say the Russian for a parcel, then the Russian for a blouse, then say \emph{futbólka} again
\end{itemize}

\begin{tcolorbox}[breakable,colback=teal!4,colframe=teal!35!black,title={Before you move on}]
What does \ru{футболка} mean? (\textquotedblleft{}A T-shirt\textquotedblright{}.) Say it once more.
\end{tcolorbox}
\section[\texorpdfstring{\ru{колготки} (kolgótki)}{колготки (kolgótki)}]{\ru{колготки} (kolgótki) — tights}

\label{lesson:RU-C162-kolgotki}



\begin{quote}
Before the new one: say the Russian for a blouse, then the Russian for a T-shirt.
\end{quote}

\subsection*{You'll want to know: \ru{колготки}}
\textbf{\ru{колготки}} (\emph{kolgótki}) — \textquotedblleft{}tights\textquotedblright{}.

\begin{cousinweb}
\textbf{\ru{колготки}} comes from Czech \emph{kalhotky}, \textquotedblleft{}panties\textquotedblright{}, the diminutive of \emph{kalhoty}, \textquotedblleft{}trousers\textquotedblright{}.
\end{cousinweb}

\begin{grammarlens}[title={What you've built}]
One more word for this chapter.
\end{grammarlens}

\subsection*{Your turn}
\begin{itemize}
\raggedright
  \item \emph{Say it:} \emph{kolgótki}
  \item \emph{Say it:} \emph{kolgótki}, once more
  \item \emph{Say it:} say \emph{kolgótki}
  \item \emph{Recall:} say the Russian for a blouse, then the Russian for a T-shirt, then say \emph{kolgótki} again
\end{itemize}

\begin{tcolorbox}[breakable,colback=teal!4,colframe=teal!35!black,title={Before you move on}]
What does \ru{колготки} mean? (\textquotedblleft{}Tights\textquotedblright{}.) Say it once more.
\end{tcolorbox}
\section[\texorpdfstring{\ru{кроссовки} (krossóvki)}{кроссовки (krossóvki)}]{\ru{кроссовки} (krossóvki) — trainers}

\label{lesson:RU-C162-krossovki}



\begin{quote}
Before the new one: say the Russian for a T-shirt, then the Russian for tights.
\end{quote}

\subsection*{You'll want to know: \ru{кроссовки}}
\textbf{\ru{кроссовки}} (\emph{krossóvki}) — \textquotedblleft{}trainers\textquotedblright{}.

\begin{cousinweb}
\textbf{\ru{кроссовки}} is built on \emph{\ru{кросс}}, \textquotedblleft{}a cross-country run\textquotedblright{}, from English \emph{cross(-country)}, with \emph{-\ru{овки}}: shoes for running.
\end{cousinweb}

\begin{grammarlens}[title={What you've built}]
One more word for this chapter.
\end{grammarlens}

\subsection*{Your turn}
\begin{itemize}
\raggedright
  \item \emph{Say it:} \emph{krossóvki}
  \item \emph{Say it:} \emph{krossóvki}, once more
  \item \emph{Say it:} say \emph{krossóvki}
  \item \emph{Recall:} say the Russian for a T-shirt, then the Russian for tights, then say \emph{krossóvki} again
\end{itemize}

\begin{tcolorbox}[breakable,colback=teal!4,colframe=teal!35!black,title={Before you move on}]
What does \ru{кроссовки} mean? (\textquotedblleft{}Trainers\textquotedblright{}.) Say it once more.
\end{tcolorbox}
\section[\texorpdfstring{\ru{тапочки} (tápochki)}{тапочки (tápochki)}]{\ru{тапочки} (tápochki) — slippers}

\label{lesson:RU-C162-tapochki}



\begin{quote}
Before the new one: say the Russian for tights, then the Russian for trainers.
\end{quote}

\subsection*{You'll want to know: \ru{тапочки}}
\textbf{\ru{тапочки}} (\emph{tápochki}) — \textquotedblleft{}slippers\textquotedblright{}.

\begin{cousinweb}
\textbf{\ru{тапочки}} is a diminutive of \emph{\ru{тапки}}, \textquotedblleft{}slippers\textquotedblright{}. Vasmer has no entry for \emph{\ru{тапки}}, and its origin is not settled.
\end{cousinweb}

\begin{grammarlens}[title={What you've built}]
One more word for this chapter.
\end{grammarlens}

\subsection*{Your turn}
\begin{itemize}
\raggedright
  \item \emph{Say it:} \emph{tápochki}
  \item \emph{Say it:} \emph{tápochki}, once more
  \item \emph{Say it:} say \emph{tápochki}
  \item \emph{Recall:} say the Russian for tights, then the Russian for trainers, then say \emph{tápochki} again
\end{itemize}

\begin{tcolorbox}[breakable,colback=teal!4,colframe=teal!35!black,title={Before you move on}]
What does \ru{тапочки} mean? (\textquotedblleft{}Slippers\textquotedblright{}.) Say it once more.
\end{tcolorbox}
